\documentclass[10pt,a4paper]{amsart}
\usepackage[margin=19mm]{geometry}
\usepackage{amsmath,amssymb,mathtools}
\usepackage[scr=rsfs,cal=euler]{mathalfa}
\usepackage{xfrac}
\usepackage[unicode,hidelinks,bookmarks=false]{hyperref}
\newcommand{\ie}{i.e.\ }
\newcommand{\cf}{cf.\ }
\newcommand{\resp}{respectively\ }
\newcommand{\Subsection}{\subsection}
\providecommand{\nobreakdash}{\nobreak\hbox{-}}
\setlength{\emergencystretch}{2em}
\multlinegap0pt
\usepackage{bm}
\usepackage{bbm}
\usepackage{dsfont}
\usepackage{xspace}
\usepackage{cancel}
\usepackage{mathtools}
\usepackage{amsthm}
\makeatletter
\@ifundefined{thm}{\newtheorem{thm}{Theorem}}{}
\@ifundefined{lem}{\newtheorem{lem}[thm]{Lemma}}{}
\@ifundefined{prop}{\newtheorem{prop}[thm]{Proposition}}{}
\makeatother
\newcommand{\Dbi}{\mathcal{D}_{\operatorname{bi}}}
\newcommand{\LieDbi}{\operatorname{Lie}\langle \Dbi\rangle}
\newcommand{\QDbi}{\QQ\langle \Dbi\rangle}
\newcommand{\QQ}{\mathbb{Q}}
\newcommand{\ariemp}{\{-,-\}\!_A}
\newcommand{\ari}[2]{\{#1,#2\}\!_A}
\newcommand{\der}[1]{d_{#1}^A}
\newcommand{\lqq}{\mathfrak{lq}}
\title[An extension of the linearized double shuffle Lie algebra]{An extension of the linearized double shuffle Lie algebra}
\author{Annika Burmester}
\date{Source: arXiv:2508.05024v1 (CC BY 4.0)}
\begin{document}
\maketitle

\noindent\emph{Source and licence.} An extension of the linearized double shuffle Lie algebra. Version arXiv:2508.05024v1 (CC BY 4.0), distributed under the Creative Commons Attribution 4.0 International licence, as recorded in the arXiv licence metadata for this exact version and shown on the abstract page https://arxiv.org/abs/2508.05024v1. Full rights notice: licenses/arxiv-2508.05024-CC-BY-4.0.txt.\par\smallskip

\noindent\textit{Editorial scope.} Counted unit: the Thm whose source label is \texttt{thm:lq\_deriv} in the article (source0.tex lines 692--693, source label \texttt{thm:lq\_deriv}), delivered together with the article's own complete original proof of it (source0.tex lines 695--741, one \texttt{proof} environment of 47 lines). No mathematics is written, repaired or re-derived: statement and proof are the article's own text line for line. Required context retained uncounted: lem:lqq\_in\_Dbi (lines 638--643); prop:d\_tau\_commute (lines 662--664). Every definition, display and label of those lines is retained unchanged. Omitted from this package and not redistributed: 1--637, 644--661, 665--691, 694--694, 742--1413. Nothing inside a retained excerpt is omitted or altered: every retained body is delivered exactly as the article's lines are written, so no omission receipt of any kind accompanies this package. Citations: the retained excerpts carry no citation command at all, so no bibliography entry of the article is transcribed and no reference is redistributed. Reference adaptation: none was needed and none was made; every reference inside the delivered unit resolves inside it, so no unresolved reference or citation is produced. Printed slips kept literal: no printed word, symbol, hypothesis or number of the counted statement or of its proof was altered. Layout and class: the article's own class is \texttt{amsart} with packages including amsmath, amsthm, amssymb, amscd, mathdots, mathtools, enumerate, shuffle, stmaryrd, faktor, hyperref, xcolor, float, scalerel; the wrapper is the lane's pinned \texttt{amsart} editorial wrapper at 10pt on A4, which restates the theorem-like environments the retained text needs and adds the article's own macro definitions that the retained text needs (\texttt{\textbackslash Dbi}, \texttt{\textbackslash LieDbi}, \texttt{\textbackslash QDbi}, \texttt{\textbackslash QQ}, \texttt{\textbackslash ari}, \texttt{\textbackslash ariemp}, \texttt{\textbackslash der}, \texttt{\textbackslash lqq}), with extra wrapper packages (bm, bbm, dsfont, xspace, cancel, mathtools). The delivered numbering is the unit's own. Assets: no retained span contains a figure environment, a graphics-inclusion command, a PDF-page inclusion, a table or a TikZ picture; no image file is copied into this package, the package declares no asset dependency and no image file is redistributed with it. Licence: version arXiv:2508.05024v1 of this article is distributed under the Creative Commons Attribution 4.0 International licence, as recorded in the arXiv licence metadata for this exact version and shown on the abstract page https://arxiv.org/abs/2508.05024v1. A full rights notice is delivered at licenses/arxiv-2508.05024-CC-BY-4.0.txt. Characters: every retained excerpt is pure ASCII, so the delivered engine, XeLaTeX with its default Latin Modern text font, reports no missing character.
\begin{lem} \label{lem:lqq_in_Dbi} The space $\lqq$ consists of all $\Phi\in \LieDbi$ such that
\begin{flalign*}
\qquad \text{(ii')} \quad & (\Phi\mid D_{k,m})=0, \qquad k+m \text{ even},\ k\geq1,\ m\geq0, &&\\
\qquad \text{(iv')} \quad & \tau_{\Dbi}(\Phi)=\Phi.
\end{flalign*}
\end{lem} 

\begin{prop} \label{prop:d_tau_commute} The maps $\tau_{\Dbi}$ and $\delta$ commute,
\[\tau_{\Dbi}\circ\delta=\delta\circ \tau_{\Dbi}.\]
\end{prop}

\begin{thm} \label{thm:lq_deriv} The tuple $(\lqq,\ariemp,\delta)$ is a differential Lie algebra.
\end{thm}

\begin{proof} Evidently, the derivation $\delta$ preserves the space $\LieDbi$ and the condition (ii') from Lemma \ref{lem:lqq_in_Dbi}. By Proposition \ref{prop:d_tau_commute}, $\delta$ also preserves the $\tau_{\Dbi}$-invariance given in in (iv'). 
We verify the compatibility of the Lie bracket $\ariemp$ and the derivation $\delta$,
\begin{align} \label{eq:lq_deriv_liebrack}
\delta\big(\ari{f}{g}\big)=\ari{\delta(f)}{g}+\ari{f}{\delta(g)},\qquad f,g\in\QDbi.
\end{align} 
By $\QQ$-linearity, we can assume that $f=D_{k_1,m_1}\cdots D_{k_d,m_d},\ g=D_{i_1,n_1}\cdots D_{i_r,n_r}$. Then, we have 
\begin{align} \label{eq:lq_deriv1}
&\delta\big(\der{f}(g)\big)= \sum_{\substack{s,t=1 \\ s\neq t}}^r D_{i_1,n_1}\cdots \der{f}(D_{i_s,n_s})\cdots \delta(D_{i_t,n_t})\cdots D_{i_r,n_r} \\
\label{eq:lq_deriv2}
&+\sum_{s=1}^r\sum_{\substack{l_1+\cdots+l_d+l=k_1+\cdots+k_d \\ p_1+\cdots+p_d+p=n_s}} (-1)^l\binom{k_1-1}{l_1-1}\cdots \binom{k_d-1}{l_d-1} \binom{n_s}{p_1,\ldots,p_d,p}  \\
&\hspace{2.4cm} D_{i_1,n_1}\cdots D_{i_{s-1},n_{s-1}}[D_{i_s+1+l,p+1},D_{l_1,m_1+p_1}\cdots D_{l_d,m_d+p_d}]D_{i_{s+1},n_{s+1}}\cdots D_{i_r,n_r}\big)  \nonumber \\
\label{eq:lq_deriv3}
&+\sum_{s=1}^r \sum_{t=1}^d\sum_{\substack{l_1+\cdots+l_d+l=k_1+\cdots+k_d \\ p_1+\cdots+p_d+p=n_s}} (-1)^l\binom{k_1-1}{l_1-1}\cdots\binom{k_d-1}{l_d-1} \binom{n_s}{p_1,\ldots,p_d,p} \\
&D_{i_1,n_1}\cdots D_{i_{s-1},n_{s-1}}[D_{i_s+l,p},D_{l_1,m_1+p_1}\cdots D_{l_t+1,m_t+1+p_1} \cdots D_{l_d,m_d+p_d}]D_{i_{s+1},n_{s+1}}\cdots D_{i_r,n_r}. \nonumber
\end{align}
For the third sum \eqref{eq:lq_deriv3}, we observe that it is equal to
\begin{align*}
&\sum_{s=1}^r \sum_{t=1}^d\sum_{\substack{l_1+\cdots+l_d+l=k_1+\cdots+k_d+1 \\ p_1+\cdots+p_d+p=n_s}} (-1)^l\binom{k_1-1}{l_1-1}\cdots \binom{k_t-1}{l_t-2}\cdots \binom{k_d-1}{l_d-1} \binom{n_s}{p_1,\ldots,p_d,p} \\
&D_{i_1,n_1}\cdots D_{i_{s-1},n_{s-1}}[D_{i_s+l,p},D_{l_1,m_1+p_1}\cdots D_{l_t,m_t+1+p_1} \cdots D_{l_d,m_d+p_d}]D_{i_{s+1},n_{s+1}}\cdots D_{i_r,n_r}\\
&=\sum_{t=1}^d\sum_{s=1}^r\sum_{\substack{l_1+\cdots+l_d+l=k_1+\cdots+k_d+1 \\ p_1+\cdots+p_d+p=n_s}} (-1)^l\binom{k_1-1}{l_1-1}\cdots \binom{k_t}{l_t-1}\cdots \binom{k_d-1}{l_d-1} \binom{n_s}{p_1,\ldots,p_d,p} \\
&D_{i_1,n_1}\cdots D_{i_{s-1},n_{s-1}}[D_{i_s+l,p},D_{l_1,m_1+p_1}\cdots D_{l_t,m_t+1+p_1} \cdots D_{l_d,m_d+p_d}]D_{i_{s+1},n_{s+1}}\cdots D_{i_r,n_r} \\
&-\sum_{s=1}^r \sum_{t=1}^d\sum_{\substack{l_1+\cdots+l_d+l=k_1+\cdots+k_d \\ p_1+\cdots+p_d+p=n_s}} (-1)^{l+1}\binom{k_1-1}{l_1-1}\cdots \binom{k_t-1}{l_t-1}\cdots \binom{k_d-1}{l_d-1} \binom{n_s}{p_1,\ldots,p_d,p} \\
&D_{i_1,n_1}\cdots D_{i_{s-1},n_{s-1}}[D_{i_s+l+1,p},D_{l_1,m_1+p_1}\cdots D_{l_t,m_t+1+p_1} \cdots D_{l_d,m_d+p_d}]D_{i_{s+1},n_{s+1}}\cdots D_{i_r,n_r} \\
&=\der{\delta(f)}(g)+\sum_{s=1}^r \sum_{t=1}^dD_{i_1,n_1}\cdots D_{i_{s-1},n_{s-1}} \der{D_{k_1,m_1}\cdots D_{k_t,m_t+1}\cdots D_{k_d,m_d}}(D_{i_s+1,n_s} )D_{i_{s+1},n_{s+1}}\cdots D_{i_r,n_r}.
\end{align*}
Next, we observe that the first and second sum \eqref{eq:lq_deriv1} and \eqref{eq:lq_deriv2} are equal to
\begin{align*}
&\sum_{\substack{s,t=1 \\ s\neq t}}^r D_{i_1,n_1}\cdots \der{f}(D_{i_s,n_s})\cdots \delta(D_{i_t,n_t})\cdots D_{i_r,n_r} \\
&+\sum_{s=1}^r\sum_{\substack{l_1+\cdots+l_d+l=k_1+\cdots+k_d \\ p_1+\cdots+p_d+p=n_s+1}} (-1)^l\binom{k_1-1}{l_1-1}\cdots \binom{k_d-1}{l_d-1} \binom{n_s}{p_1,\ldots,p_d,p-1}  \\
&\hspace{2.4cm} D_{i_1,n_1}\cdots D_{i_{s-1},n_{s-1}}[D_{i_s+1+l,p},D_{l_1,m_1+p_1}\cdots D_{l_d,m_d+p_d}]D_{i_{s+1},n_{s+1}}\cdots D_{i_r,n_r}\big)\\
&=\sum_{\substack{s,t=1 \\ s\neq t}}^r D_{i_1,n_1}\cdots \der{f}(D_{i_s,n_s})\cdots \delta(D_{i_t,n_t})\cdots D_{i_r,n_r} \\
&+\sum_{s=1}^r\sum_{\substack{l_1+\cdots+l_d+l=k_1+\cdots+k_d \\ p_1+\cdots+p_d+p=n_s+1}} (-1)^l\binom{k_1-1}{l_1-1}\cdots \binom{k_d-1}{l_d-1} \binom{n_s+1}{p_1,\ldots,p_d,p}  \\
&\hspace{2.4cm} D_{i_1,n_1}\cdots D_{i_{s-1},n_{s-1}}[D_{i_s+1+l,p},D_{l_1,m_1+p_1}\cdots D_{l_d,m_d+p_d}]D_{i_{s+1},n_{s+1}}\cdots D_{i_r,n_r}\big)\\
&-\sum_{s=1}^r \sum_{t=1}^d\sum_{\substack{l_1+\cdots+l_d+l=k_1+\cdots+k_d \\ p_1+\cdots+p_d+p=n_s}} (-1)^l\binom{k_1-1}{l_1-1}\cdots \binom{k_d-1}{l_d-1} \binom{n_s}{p_1,\ldots,p_d,p}\\ 
&\hspace{1.4cm} D_{i_1,n_1}\cdots D_{i_{s-1},n_{s-1}}[D_{i_s+1+l,p},D_{l_1,m_1+p_1}\cdots D_{l_t,m_t+1+p_t}\cdots D_{l_d,m_d+p_d}]D_{i_{s+1},n_{s+1}}\cdots D_{i_r,n_r}\big)\\
&=\der{f}(\delta(g))-\sum_{s=1}^r \sum_{t=1}^d D_{i_1,n_1}\cdots D_{i_{s-1},n_{s-1}}\der{D_{k_1,m_1}\cdots D_{k_t,m_t+1}\cdots D_{k_d,m_d}}(D_{i_s+1,n_s})D_{i_{s+1},n_{s+1}}\cdots D_{i_r,n_r}.
\end{align*}
Both calculations together show that
\begin{align} \label{eq:lq_deriv_deriv}
\delta\big(\der{f}(g)\big)=\der{\delta(f)}(g)+\der{f}(\delta(g)).
\end{align}
Since $\delta$ is a derivation, we have
\begin{align} \label{eq:lq_deriv_comm}
\delta\big([f,g]\big)=[\delta(f),g]+[f,\delta(g)].
\end{align}
From \eqref{eq:lq_deriv_deriv} and \eqref{eq:lq_deriv_comm}, we derive the claimed identity \eqref{eq:lq_deriv_liebrack}.
\end{proof}

\end{document}
